\scanpage{78}
\mspage{II. 26}
% archive sheet number: 77
% lines 1–7
connexe, \inserted{\unsure{i.e.} son image inverse $\bar C_1$ sur $C$ est géom. connexe, ou même géom. irréduct.} Or $X_1$ provient d'un revêtement \deleted{étale} \inserted{normal} intègre \deleted{\unsure{ramifié}} \inserted{[EGA \unsure{V|IV}]} $X'_1-X'$. En vertu du th. de Bertini, l'image inverse \deleted{\unsure{de}} dans $X'_1\longrightarrow\mathbb P^r$ \deleted{\unsure{est géom}} de $L_t$ est \deleted{géom.} géom. irréductible, ce qui implique que l'ouvert $C_1^{\unsure{'|-}}$ tracé sur $C$ l'est également, i.e. $\bar C_1$ est géom connexe. Cela prouve le lemme.
% lines 8–14
3.3. De 3.2 résulte aussitôt que si \deleted{l'extension} $C$ sur $K$ satisfait aux conclusions de 1.1, il en est de même de $X$ sur $k$. Compte tenu de la réduction déjà faite dans 2.3, cela nous ramène, pour prouver 1.1 \inserted{sur $k$}, au cas d'une courbe algébrique \inserted{sur $K$}, déjà traité au n° 2.

(NB. Si $k$ est de t.f., il en est de même de $K$ !).
% lines 15–23
\uline{Remarque} 3.4. Le lemme 3.2. permet aussi de montrer que si $X$ est un schéma de t.f. sur un corps séparablement clos \inserted{de car. $p>0$}, alors le \deleted{\unsure{petit}} plus grand \inserted{groupe} quotient profini premier à $p$ de $\pi_1(X,\xi)$, \inserted{soit $\pi_1(X,\xi)'$} est topologique à engendrement fini. En effet, la théorie de la descente \inserted{(pour les groupes fond$^{\unsure{aux}}$ (SGA 1 \unsure{IX|VIII}))} nous permet de nous ramener au cas où $X$ est géom. unibranche irréductible, \deleted{\unsure{et}} puis au cas \deleted{\unsure{d'un}} $X$ affine, puis 3.2. nous ramène au cas où $X$ \deleted{\unsure{d'une}} \inserted{est un ouvert} d'une courbe projective lisse $X'$. Or pour un tel $X$, cela résulte de la structure connue de $\pi_1(X,\xi)'$ (SGA 1\ldots).
